\section{附录}

\subsection{关键代码示例}

\subsubsection{Claude Code 风格 TAOR 主循环}

\begin{lstlisting}[language=python, caption=apps/api/agent/agent\_loop.py(节选)]
class AgentLoop:
    async def run(self, initial_state, system_prompt, can_use_tool):
        state = initial_state
        try:
            while True:
                # 1) 预算与轮数检查
                if state.turn_count > self.max_turns:
                    yield Terminal(reason=TerminalReason.MAX_TURNS, ...)
                    return

                # 2) 上下文压缩
                messages = await self.compression.run_full_pipeline(
                    state.messages, state.auto_compact_tracking
                )

                # 3) 调用 LLM
                response = await self.llm_call(messages=messages, ...)
                tool_calls = response.get("tool_calls", [])

                # 4) 没 tool calls -> Stop hook 链
                if not tool_calls:
                    stop_result = await self.hooks.execute_stop({...})
                    if stop_result.prevent_continuation:
                        yield Terminal(reason=TerminalReason.STOP_HOOK_PREVENTED)
                        return
                    yield Terminal(reason=TerminalReason.COMPLETED, ...)
                    return

                # 5) 执行工具(经过 PreToolUse hook)
                tool_results = await self._execute_tools(tool_calls, state)

                # 6) 整对象重建 State
                state = replace(
                    state,
                    messages=state.messages + [response.message] + tool_results,
                    turn_count=state.turn_count + 1,
                    transition=Transition(reason=ContinueReason.NEXT_TURN),
                )
        except asyncio.CancelledError:
            yield Terminal(reason=TerminalReason.ABORTED_STREAMING, ...)
\end{lstlisting}

\subsubsection{Oracle 早期信号雷达 - 12 数据源并行}

\begin{lstlisting}[language=python, caption=apps/api/agents/oracle\_agent.py(节选)]
class OracleAgent:
    def __init__(self):
        self.sources = [
            PolymarketSource(),    # 国际预测市场
            ManifoldSource(),      # 社区预测
            HackerNewsSource(),    # 科技叙事
            GitHubTrendingSource(),# 开源讨论
            WeiboMockSource(),     # 微博实时
            WeixinIndexMockSource(),# 微信指数
            BaiduIndexMockSource(),# 百度指数
            XHSInspirationMockSource(),# 小红书灵感
            JuliangTrendMockSource(),  # 抖音算数
        ]

    async def scan(self, topic_seed, category, target_platforms):
        # 12 数据源并行
        results = await asyncio.gather(
            *[s.fetch(topic_seed) for s in self.sources],
            return_exceptions=True,
        )

        # CUSUM + MAD-zscore 异常检测
        all_samples = []
        for r in results:
            if isinstance(r, list):
                all_samples.extend(r)

        # 矛盾推理(借鉴 Digital Oracle)
        contradictions = find_contradictions(all_samples)

        # 聚合成 Top 10 趋势 + 置信度
        trends = self._aggregate_trends(all_samples, category, target_platforms)

        return OracleReport(trends=trends, ...)
\end{lstlisting}

\subsubsection{多 Agent 论坛辩论}

\begin{lstlisting}[language=python, caption=apps/api/agents/forum\_debate\_agent.py(节选)]
class ForumDebateAgent:
    """1 主持 + 3 专家(数据派/创意派/风险派)辩论 2 轮"""
    
    async def debate(self, topic, context, max_rounds=2):
        # Round 1: 各自独立发言(并行)
        responses_r1 = await asyncio.gather(*[
            self._speaker_speak(role, prompt)
            for role in ["data_expert", "creative_expert", "risk_expert"]
        ])

        # Round 2-N: 互相反驳
        for r in range(2, max_rounds + 1):
            responses_rn = await asyncio.gather(*[
                self._speaker_speak(role, build_rebuttal_prompt(...))
                for role in speakers
            ])

        # 主持人裁判
        decision = await self._host_decide(topic, all_views)

        return DebateResult(rounds=rounds, final_decision=decision, ...)
\end{lstlisting}

% --------------------------------------------

\subsection{数据来源与文献}

\subsubsection{学术论文}

\begin{itemize}
\item PopSim: arxiv:2512.02533 — LLM 多智能体社交沙箱传播预测
\item AutoCas: arxiv:2502.18040 — 级联扩散自回归 LLM 模型
\item BuzzProphet: 微博 hashtag 流行度预测(2024-2025)
\item HyperFusion: SMP Challenge 2025 多模态融合
\item Loewenstein 1994: Curiosity Gap 理论(Psychological Bulletin)
\end{itemize}

\subsubsection{行业报告}

\begin{itemize}
\item 克劳锐《2024 年 KOL 发展年报》:中国万粉以上创作者 1500 万+
\item Goldman Sachs Research:全球 Creator Economy 2027 \$4800 亿
\item Ebiquity:中国 2025 KOL 营销市场 ¥930 亿
\item 36氪《全中国最懂小红书的「郑州帮」》:KOC 真实日常
\item 36氪《MCN 不是网红们的救世主》:MCN-达人关系矛盾
\end{itemize}

\subsubsection{开源项目(灵感来源)}

\begin{itemize}
\item Digital Oracle: \url{https://github.com/komako-workshop/digital-oracle}
\item Claude Code 51 万行: 基于第三方逆向 \texttt{alejandrobalderas/claude-code-from-source}
\item MiroFish: \url{https://github.com/666ghj/MiroFish}
\item BettaFish: \url{https://github.com/666ghj/BettaFish}
\item Shannon: \url{https://github.com/Kocoro-lab/Shannon}
\item OASIS: \url{https://github.com/camel-ai/oasis}
\end{itemize}

\subsubsection{法规与标准}

\begin{itemize}
\item 《人工智能生成合成内容标识办法》(2025-09-01 施行)
\item 《互联网信息服务深度合成管理规定》
\item 《生成式人工智能服务管理暂行办法》
\item 《个人信息保护法》(PIPL)
\item OWASP GenAI / LLM Top 10 v2 (2024-11)
\end{itemize}

\subsubsection{腾讯生态文档}

\begin{itemize}
\item 腾讯混元 OpenAI 兼容: \url{https://cloud.tencent.com/document/product/1729/111007}
\item 元器 Agent API: \url{https://yuanqi.tencent.com/guide/publish-agent-api-documentation}
\item 视频号助手 API: \url{https://developers.weixin.qq.com/doc/channels/api/channels/}
\item 微信公众平台开发文档: \url{https://developers.weixin.qq.com}
\end{itemize}

% --------------------------------------------

\subsection{致谢}

感谢腾讯 PCG 校园 AI 产品创意大赛组委会提供这次机会。

感谢以下开源项目与作者(灵感谱系):

\begin{itemize}
\item \textbf{komako-workshop/digital-oracle}: 早期信号思想的奠基性工作
\item \textbf{666ghj/MiroFish + BettaFish + MindSpider}: 群体仿真与多 Agent 辩论
\item \textbf{Anthropic Claude Code}: 51 万行代码架构纪律的标杆
\item \textbf{Kocoro-lab/Shannon}: 生产级多 Agent 编排
\item \textbf{NousResearch/Hermes Agent}: 记忆系统设计
\item \textbf{LangGraph / LiteLLM / FastAPI / Tauri / Streamlit}: 工程基石
\end{itemize}

感谢 36氪、克劳锐、CSDN 上分享真实工作日常的 KOC 们,
你们的「不可言说的痛点」是 Ripple 设计的起点。

\bigskip

\begin{flushright}
\textit{Ripple Team \\
2026 年 4 月 28 日}
\end{flushright}
